\documentclass[../../main.tex]{subfiles}

\begin{document}

\section{Radon 测度的构造}\label{section:ch21-radon构造}

令 $X$ 为紧 Hausdorff 空间. 我们希望构造 $\mathcal{B}(X)$ 上的 Radon 测度. 我们首先定义开集族 $\mathcal{T}$ 上的集函数 $\lambda:\mathcal{T}\to[0,\infty]$ 并证明 $\lambda$ 是预测度. 从 Carath\'eodory-Hahn 定理推出由 $\lambda$ 诱导的外测度 $\lambda^{*}$ 在包含 $\mathcal{B}(X)$ 的 $\sigma$ 代数 $\mathcal{M}$ 上诱导了一个测度 $\overline{\lambda}$. 然而, 该定理需要 $\lambda$ 在 $\mathcal{T}$ 上可数可加. 一般地这是不正确的. 因此我们从 $\mathcal{T}$ 上的集函数 $\lambda:\mathcal{T}\to[0,\infty]$ 开始, 它具有以下两个性质:
\begin{enumerate}[(i)]
\item 单调性: 若 $\mathcal{O}_{1}$ 和 $\mathcal{O}_{2}$ 是开的且 $\mathcal{O}_{1}\subseteq\mathcal{O}_{2}$, 则 $\lambda(\mathcal{O}_{1})\leqslant\lambda(\mathcal{O}_{2})$.
\item 次可加性: 对开集的可数族 $\{\mathcal{O}_{k}\}_{k=1}^{\infty}$, $\lambda\left(\bigcup_{k=1}^{\infty}\mathcal{O}_{k}\right)\leqslant\sum_{k=1}^{\infty}\lambda(\mathcal{O}_{k})$.
\end{enumerate}

这些性质保证 $\lambda^{*}$ 是外测度. 我们定义 $\lambda^{*}:2^{X}\to[0,\infty]$ 为
\begin{align*}
\text{对 }E\subseteq X,\quad\lambda^{*}(E)=\inf\{\lambda(\mathcal{O})\mid E\subseteq\mathcal{O},\ \mathcal{O}\text{ 是开的}\}.
\end{align*}
从 $\lambda$ 的单调性和次可加性推出 $\lambda^{*}$ 是 $2^{X}$ 上的外测度. 此外, 若 $\lambda$ 有进一步的性质, 则 $\lambda^{*}$ 是 Carath\'eodory 外测度, 因此 Borel 集是 $\lambda^{*}$ 可测的. 这样的性质是什么呢? 我们需要以下关于 $\lambda$ 的假设: 对任何两个开集 $\mathcal{O}_{1}$ 和 $\mathcal{O}_{2}$ 使得 $\overline{\mathcal{O}}_{1}\cap\overline{\mathcal{O}}_{2}=\varnothing$,
\begin{align}\label{align:ch21-不交开集可加}
\lambda(\mathcal{O}_{1}\cup\mathcal{O}_{2})=\lambda(\mathcal{O}_{1})+\lambda(\mathcal{O}_{2}).
\end{align}
从这个性质推出, 若 $K_{1}$ 和 $K_{2}$ 是紧集且 $K_{1}\cap K_{2}=\varnothing$, 则
\begin{align*}
\lambda^{*}(K_{1}\cup K_{2})=\lambda^{*}(K_{1})+\lambda^{*}(K_{2}).
\end{align*}
这意味着 $\lambda^{*}$ 是 Carath\'eodory 外测度. 因此 Borel 集是 $\lambda^{*}$ 可测的.

现在定义 $\mathcal{M}$ 为 $\lambda^{*}$ 可测集的 $\sigma$ 代数. 测度 $\mu=\lambda^{*}|_{\mathcal{B}(X)}$ 是 Borel 测度. 我们需要证明它是 Radon 测度, 即每个 Borel 集是正则的. 这通过以下步骤完成:
\begin{enumerate}[(1)]
\item 每个开集 $\mathcal{O}$ 是外正则的, 这从 $\lambda^{*}$ 的定义立即得出.
\item 每个紧集 $K$ 是内正则的, 这从 $\lambda^{*}(K)=\inf\{\lambda(\mathcal{O})\mid K\subseteq\mathcal{O}\}$ 与紧集的分离性质得出.
\item 从 Borel 集的正则性和测度的连续性质推出一般 Borel 集的正则性.
\end{enumerate}

因此我们有以下定理.

\begin{theorem}[Radon 测度构造定理]\label{theorem:ch21-radon构造}
令 $X$ 为紧 Hausdorff 空间而 $\lambda:\mathcal{T}\to[0,\infty]$ 是定义在 $X$ 的开子集族上的集函数, 它满足以下性质:
\begin{enumerate}[(i)]
\item 对开集 $\mathcal{O}_{1}$ 和 $\mathcal{O}_{2}$, 若 $\mathcal{O}_{1}\subseteq\mathcal{O}_{2}$, 则 $\lambda(\mathcal{O}_{1})\leqslant\lambda(\mathcal{O}_{2})$.
\item 对开集的可数族 $\{\mathcal{O}_{k}\}_{k=1}^{\infty}$, $\lambda\left(\bigcup_{k=1}^{\infty}\mathcal{O}_{k}\right)\leqslant\sum_{k=1}^{\infty}\lambda(\mathcal{O}_{k})$.
\item 对具有不相交闭包的开集 $\mathcal{O}_{1}$ 和 $\mathcal{O}_{2}$, $\lambda(\mathcal{O}_{1}\cup\mathcal{O}_{2})=\lambda(\mathcal{O}_{1})+\lambda(\mathcal{O}_{2})$.
\end{enumerate}
则存在 $\mathcal{B}(X)$ 上的唯一 Radon 测度 $\mu$ 使得对每个开集 $\mathcal{O}$,
\begin{align*}
\mu(\mathcal{O})=\lambda(\mathcal{O}).
\end{align*}
\end{theorem}

\subsection*{习题}

\begin{exercise}\label{exercise:ch21-25}
证明由 $\lambda$ 通过\eqref{align:ch21-不交开集可加}诱导的外测度 $\lambda^{*}$ 是 Carath\'eodory 外测度.
\end{exercise}

\begin{exercise}\label{exercise:ch21-26}
证明每个紧集关于 $\mu$ 是内正则的.
\end{exercise}

\begin{exercise}\label{exercise:ch21-27}
证明每个 Borel 集关于 $\mu$ 是正则的. (提示: 首先证明具有有限外测度的 Borel 集的族是 $\sigma$ 代数.)
\end{exercise}

\end{document}